\documentclass[12pt]{article}
\usepackage{amsfonts,amsthm,amsmath,amssymb,mathtools}
\usepackage{mathrsfs}
\usepackage{enumitem}
\usepackage{tikz-cd}
\usepackage{geometry}
\usepackage{mlmodern}

\geometry{letterpaper, portrait, margin=1in}

\setcounter{section}{0}

\renewcommand{\thesection}{\S\arabic{section}}
\renewcommand{\thesubsection}{\arabic{section}}
\newcommand{\x}{\mathbf{x}}
\DeclareMathOperator{\diam}{diam}
\newcommand{\dd}{\mathop{}\!\mathrm{d}}

\newtheorem{thm}{Theorem}
\newenvironment{theorem}[1]{
	\renewcommand{\thethm}{#1}
	\begin{thm}
		}{\end{thm}}

\newtheorem{lma}{Lemma}
\newenvironment{lemma}[1]{
	\renewcommand{\thelma}{#1}
	\begin{lma}
		}{\end{lma}}

\theoremstyle{definition}
\newtheorem{ex}{}
\newenvironment{exercise}[1]{
	\renewcommand{\theex}{#1}
	\begin{ex}
		}{\end{ex}}

\title{Homework 12}
\author{Connor Mooneyhan}
\date{April 23, 2026}

\begin{document}

\maketitle

\begin{ex}
	Let $(e_j)$ be an orthonormal sequence in a Hilbert space $H$.
	Show that if $x = \sum_{j=1}^\infty a_je_j$ and $y = \sum_{j=1}^\infty b_je_j$, then \begin{equation*}
		\langle x, y \rangle = \sum_{j=1}^\infty a_j\overline{b_j}.
	\end{equation*}
\end{ex}
\begin{proof}
	We have \begin{align*}
		\langle x, y \rangle & = \left\langle \sum_{i=1}^\infty a_ie_i, \sum_{j=1}^\infty b_je_j \right\rangle             \\
		                     & = \sum_{i=1}^\infty a_i \left\langle e_i, \sum_{j=1}^\infty b_je_j \right\rangle            \\
		                     & = \sum_{i=1}^\infty \sum_{j=1}^\infty a_i\overline{b_j} \left\langle e_i, e_j \right\rangle \\
		                     & = \sum_{i=1}^\infty a_i\overline{b_i} \left\langle e_i, e_i \right\rangle                   \\
		                     & = \sum_{i=1}^\infty a_i\overline{b_i}
	\end{align*}
	as desired, where the fourth equality comes from the orthogonality of $e_i$ with $e_j$ when $i \neq j$, and the final equality comes from the normality of each $e_i$.
\end{proof}

\begin{ex}
	Let $(e_j)$ be an orthonormal sequence in a Hilbert space $H$, and let $U = \mathrm{span}\lbrace e_1, e_2, \dots \rbrace$.
	Show that for any $x \in H$, the vector $y = \sum_{j=1}^\infty \langle x, e_j\rangle e_j$ satisfies $y \in \overline{U}$ and $x - y \in (\overline{U})^\perp$.
\end{ex}
\begin{proof}
	Since $U$ is a subspace, we can use the fact that $U^{\perp\perp}$ is the smallest subspace containing $U$ (proved in a previous homework), and so $U^{\perp\perp} = \overline{U}$.
	In fact, since $U^\perp$ is closed, we have $\overline{U}^\perp = U^{\perp\perp\perp} = U^\perp$.
	We thus begin by showing that $x - y \in \overline{U}^\perp$ by showing that $x - y \in U^\perp$.

	Indeed, we take the inner product of $x - y$ with a general member of $U$ to get \begin{align*}
		\left\langle x - y, \sum_{j=1}^\infty a_je_j \right\rangle & = \left\langle x - \sum_{i=1}^\infty \langle x, e_i \rangle e_i, \sum_{j=1}^\infty a_je_j \right\rangle                                                                              \\
		                                                           & = \left\langle x, \sum_{j=1}^\infty a_je_j \right\rangle - \left\langle \sum_{i=1}^\infty \langle x, e_i \rangle e_i, \sum_{j=1}^\infty a_je_j \right\rangle                         \\
		                                                           & = \sum_{j=1}^\infty \overline{a_j} \left\langle x, e_j \right\rangle - \sum_{i=1}^\infty \sum_{j=1}^\infty \langle x, e_i \rangle \overline{a_j} \left\langle e_i, e_j \right\rangle \\
		                                                           & = \sum_{j=1}^\infty \overline{a_j} \left\langle x, e_j \right\rangle - \sum_{i=1}^\infty \langle x, e_i \rangle \overline{a_i} \left\langle e_i, e_i \right\rangle                   \\
		                                                           & = \sum_{j=1}^\infty \overline{a_j} \left\langle x, e_j \right\rangle - \sum_{i=1}^\infty \langle x, e_i \rangle \overline{a_i}                                                       \\
		                                                           & = 0,
	\end{align*}
	where the fourth equality comes from the orthogonality of $e_i$ and $e_j$ when $i \neq j$, and the fifth equality comes from the normality of each $e_i$.
	Since our choice of element from $U$ was arbitrary, $x - y \in U^\perp = \overline{U}^\perp$.
\end{proof}

\begin{ex}
	Prove the following theorem.
	\\
	\textbf{Theorem 1.} Let $H$ be a Hilbert space, and $M \subset H$ be an orthonormal set.
	Then the following statements are equivalent: \begin{enumerate}[label=(\arabic*)]
		\item $M$ is total.
		\item The only $x \in H$ that is orthogonal to all vectors in $M$ is 0.
		\item $\sum_{e \in M} |\langle x, e \rangle |^2 = ||x||^2$ for all $x \in H$.
	\end{enumerate}
\end{ex}
\begin{proof}
	($(1) \implies (2)$) If $M$ is total, then $H = \overline{\mathrm{span}(M)} = \mathrm{span}(M)^{\perp\perp}$.
	Also, if $y \in M^\perp$, then $y \in \mathrm{span}(M)^\perp$ since \begin{align*}
		\left\langle y, \sum_{m \in M} a_m m \right\rangle = \sum_{m \in M} \overline{a_m}\left\langle y, m  \right\rangle = 0
	\end{align*}
	since $y$ is orthogonal to each $m$, noting that we are considering such linear combinations when $a_m = 0$ for all but finitely many $m$, since that set is exactly the span of $M$.
	In other words, given some $y \in M^\perp$, every possible choice of $x \in H$ satisfies $\langle x, y \rangle = 0$ when $y$ is orthogonal to all vectors in $M$.
	Thus $y$ is orthogonal to every vector in $H$, and the only vector which has that property is 0, so $y = 0$.

	($(2) \implies (1)$) Since the set $M^\perp = \lbrace 0 \rbrace$ is closed, we have that \begin{align*}
		H & = M^\perp \oplus M^{\perp\perp}                        \\
		  & = M^\perp \oplus \overline{\mathrm{span}(M)}           \\
		  & = \lbrace 0 \rbrace \oplus \overline{\mathrm{span}(M)} \\
		  & = \overline{\mathrm{span}(M)},
	\end{align*}
	so $M$ is total as desired.

	($(1) \implies (3)$) Assume $M$ is total.
	If $x \in \mathrm{span}(M)$, then for some finite $F \subset M$, \begin{align*}
		||x||^2 & = \left|\left| \sum_{e \in F} \langle x, e \rangle e \right|\right|^2 \\
            & = \left|\left| \sum_{e \in M} \langle x, e \rangle e \right|\right|^2 \hspace{2em}\text{(since $x$ is orthogonal to every $e \in M \setminus F$)} \\
            & = \left\langle \sum_{e \in M} \langle x, e \rangle e, \sum_{f \in M} \langle x, f \rangle f \right\rangle \\
            & = \sum_{e \in M} \sum_{f \in M}\langle x, e \rangle \overline{\langle x, f \rangle}\left\langle  e,  f \right\rangle \\
            & = \sum_{e \in M} \langle x, e \rangle \overline{\langle x, e \rangle}\left\langle  e,  e \right\rangle \hspace{2em}\text{(since $e \perp f$ when $e \neq f$)} \\
            & = \sum_{e \in M} \langle x, e \rangle \overline{\langle x, e \rangle} \\
            & = \sum_{e \in M} |\langle x, e \rangle|^2
	\end{align*}
  as desired.

  Now if $x \in H \setminus \mathrm{span}(M) = \overline{\mathrm{span}(M)} \setminus \mathrm{span}(M)$, then

  ($(3) \implies (1)$) 
\end{proof}

\end{document}
